\documentclass[11pt]{article}

\usepackage[a4paper,margin=2.5cm]{geometry}
\usepackage{amsmath,amssymb,amsthm,mathtools}
\usepackage{enumitem}
\usepackage{hyperref}

\title{Notes on the \(h_1\)-window, \(h_1\)-eigenfunction problem}
\author{}
\date{}

\newtheorem{theorem}{Theorem}[section]
\newtheorem{proposition}[theorem]{Proposition}
\newtheorem{lemma}[theorem]{Lemma}
\newtheorem{remark}[theorem]{Remark}
\newtheorem{definition}[theorem]{Definition}

\newcommand{\C}{\mathbb C}
\newcommand{\R}{\mathbb R}
\newcommand{\1}{\mathbf 1}
\newcommand{\dd}{\,d}
\newcommand{\pw}{\partial_w}
\newcommand{\pwb}{\partial_{\bar w}}

\begin{document}

\maketitle

\section{Purpose of this note}

The goal of this note is to record what is currently justified, and what is not, in the
\((h_1,h_1,U)\) localization problem.

More precisely, we consider a bounded domain \(U\subset \C\) and assume that \(h_1\) is an
eigenfunction of the \(h_1\)-window localization operator:
\[
\mathcal L_U^{h_1}h_1=\lambda h_1
\]
for some \(\lambda\in \R\).

An earlier version of this note claimed that the Bargmann identity leads, through a
Paley--Wiener / Ehrenpreis--Malgrange argument, to a compactly supported \emph{Poisson}
equation. That step is not correct as stated. The data furnished by the Bargmann identity
naturally lead instead to a compactly supported \emph{first-order} equation for
\(\partial_w u\). If one differentiates that first-order equation, one does \emph{not}
recover the clean Poisson equation previously used in the free-boundary part of the argument:
extra first-order and second-order boundary distributions appear.

Accordingly, this note does not prove the disk theorem. Instead, it organizes the corrected
operator-level conclusions:
\begin{enumerate}[label=\textnormal{(\arabic*)}]
\item the Bargmann formulation of the \((h_1,h_1,U)\) problem;
\item the correct compactly supported first-order distributional equation;
\item the exact second-order equation obtained after applying \(\partial_{\bar w}\);
\item the reason the previous Laplacian-based rigidity argument does not presently go through.
\end{enumerate}

\section{Bargmann formulation}

We assume throughout that \(U\subset \C\) is bounded. Under the Bargmann transform,
\[
h_1 \longleftrightarrow e_1(w)=\sqrt{\pi}\,w.
\]
If \(T_z\) denotes the Weyl operator on Fock space, then
\[
(T_zF)(w)=e^{\pi \bar z w-\frac{\pi}{2}|z|^2}F(w-z),
\]
and therefore
\[
T_ze_1(w)=\sqrt{\pi}\,e^{\pi \bar z w-\frac{\pi}{2}|z|^2}(w-z).
\]

Using the orthonormal basis \(e_k(w)=\sqrt{\pi^k/k!}\,w^k\), one computes
\[
\langle e_1,T_ze_1\rangle_{\mathcal F^2}
=
e^{-\frac{\pi}{2}|z|^2}(1-\pi|z|^2).
\]
Consequently, the eigenvalue equation is equivalent to
\begin{equation}\label{eq:bargmann-h1}
\int_U (1-\pi|z|^2)(w-z)e^{\pi w\bar z}e^{-\pi|z|^2}\,dA(z)=\lambda w,
\qquad w\in \C.
\end{equation}

Taking complex conjugates and replacing \(\bar w\) by \(w\), we obtain the equivalent identity
\begin{equation}\label{eq:bargmann-h1-conjugated}
\int_U (1-\pi|z|^2)(w-\bar z)e^{\pi wz}e^{-\pi|z|^2}\,dA(z)=\lambda w,
\qquad w\in \C.
\end{equation}

\section{The corrected compact-support equation}

Set
\[
g(w):=(1-\pi|w|^2)e^{-\pi|w|^2}.
\]
Define the compactly supported distribution
\begin{equation}\label{eq:def-omega}
\omega:=
\left(-\frac{1}{\pi}\pw-\bar w\right)\!\left(g(w)\1_U\right)
+\frac{\lambda}{\pi}\,\pw\delta_0.
\end{equation}

\begin{proposition}\label{prop:first-order}
For every \(w\in \C\),
\[
\langle \omega,e^{\pi wz}\rangle = 0.
\]
Equivalently,
\[
\omega
=
\bar w\,e^{-\pi|w|^2}\1_U
-\frac{1}{\pi}(1-\pi|w|^2)e^{-\pi|w|^2}\pw\1_U
+\frac{\lambda}{\pi}\,\pw\delta_0.
\]
\end{proposition}

\begin{proof}
Pairing \(\omega\) against \(z\mapsto e^{\pi wz}\), we get
\[
\left\langle
\left(-\frac{1}{\pi}\pw-\bar z\right)\!\left(g(z)\1_U\right),
e^{\pi wz}
\right\rangle
=
w\int_U g(z)e^{\pi wz}\,dA(z)
-\int_U \bar z\,g(z)e^{\pi wz}\,dA(z),
\]
because \(\pw e^{\pi wz}=\pi w e^{\pi wz}\). Also,
\[
\frac{\lambda}{\pi}\langle \pw\delta_0,e^{\pi wz}\rangle
=
-\frac{\lambda}{\pi}\cdot \pi w
=
-\lambda w.
\]
Hence
\[
\langle \omega,e^{\pi wz}\rangle
=
\int_U g(z)(w-\bar z)e^{\pi wz}\,dA(z)-\lambda w,
\]
which vanishes by \eqref{eq:bargmann-h1-conjugated}.

The expanded form follows from the identity
\[
-\frac{1}{\pi}\pw g(w)-\bar w\,g(w)=\bar w\,e^{-\pi|w|^2}.
\]
\end{proof}

\begin{remark}\label{rem:ME}
The vanishing relation in Proposition~\ref{prop:first-order} is naturally adapted to the
first-order operator \(\pw\), not to the Laplacian. Thus the honest output of the
Paley--Wiener / Ehrenpreis--Malgrange step is the existence of a compactly supported
distribution \(u\) such that
\begin{equation}\label{eq:first-order-u}
\pw u = \omega.
\end{equation}
The earlier claim that the same input yields a compactly supported solution of
\(\Delta u=\omega\) is not justified by this argument.
\end{remark}

\section{The differentiated equation}

Assume from now on that \(u\) is a compactly supported distribution satisfying
\eqref{eq:first-order-u}. Applying \(\pwb\) to both sides gives
\[
\pwb\pw u = \pwb \omega.
\]
Since \(\Delta=4\pw\pwb=4\pwb\pw\), this becomes
\[
\frac14 \Delta u = \pwb\omega.
\]

\begin{proposition}\label{prop:second-order}
In the distributional sense,
\begin{align}
\frac14 \Delta u
&=
(1-\pi|w|^2)e^{-\pi|w|^2}\1_U
+\bar w e^{-\pi|w|^2}\pwb \1_U \notag\\
&\quad
+w(2-\pi|w|^2)e^{-\pi|w|^2}\pw \1_U
-\frac1\pi(1-\pi|w|^2)e^{-\pi|w|^2}\pwb\pw \1_U
+\frac{\lambda}{4\pi}\Delta\delta_0. \label{eq:second-order}
\end{align}
Equivalently,
\begin{align}
\Delta u
&=
4(1-\pi|w|^2)e^{-\pi|w|^2}\1_U
+4\bar w e^{-\pi|w|^2}\pwb \1_U \notag\\
&\quad
+4w(2-\pi|w|^2)e^{-\pi|w|^2}\pw \1_U
-\frac4\pi(1-\pi|w|^2)e^{-\pi|w|^2}\pwb\pw \1_U
+\frac{\lambda}{\pi}\Delta\delta_0. \label{eq:second-order-full}
\end{align}
\end{proposition}

\begin{proof}
Starting from \eqref{eq:def-omega},
\[
\frac14\Delta u
=
\pwb\omega
=
\pwb\left[\left(-\frac{1}{\pi}\pw-\bar w\right)(g\1_U)\right]
+\frac{\lambda}{\pi}\pwb\pw \delta_0.
\]
Expanding gives
\[
\frac14\Delta u
=
-\frac{1}{\pi}\pwb\pw(g\1_U)
-\pwb(\bar w g\1_U)
+\frac{\lambda}{4\pi}\Delta\delta_0.
\]
Now
\[
\pw g=-\pi \bar w(2-\pi|w|^2)e^{-\pi|w|^2},
\qquad
\pwb g=-\pi w(2-\pi|w|^2)e^{-\pi|w|^2},
\]
and a direct computation yields
\[
-\frac1\pi\pwb\pw g-g-\bar w\,\pwb g
=(1-\pi|w|^2)e^{-\pi|w|^2}.
\]
Substituting these identities into the product-rule expansion proves
\eqref{eq:second-order}, and multiplying by \(4\) gives \eqref{eq:second-order-full}.
\end{proof}

\begin{lemma}\label{lem:boundary-decomposition}
Assume \(\partial U\) is of class \(C^1\), and let \(\nu=(\nu_1,\nu_2)\) denote the outward
unit normal and \(\tau\) the positively oriented unit tangent along \(\partial U\). Then
\[
\pw \1_U
=
-\frac12(\nu_1-i\nu_2)\,\delta_{\partial U},
\qquad
\pwb \1_U
=
-\frac12(\nu_1+i\nu_2)\,\delta_{\partial U}.
\]
Moreover,
\[
\langle \Delta \1_U,\varphi\rangle
=
\int_{\partial U}\partial_\nu \varphi\,ds
\qquad
(\varphi\in C_c^\infty(\C)),
\]
so \(\Delta \1_U\) is a derivative-type boundary distribution.

If we write
\[
g(w):=(1-\pi|w|^2)e^{-\pi|w|^2},
\]
then \eqref{eq:second-order-full} becomes
\begin{align}
\Delta u
&=
4g(w)\1_U
-2e^{-\pi|w|^2}
\Big(
\bar w(\nu_1+i\nu_2)
+w(2-\pi|w|^2)(\nu_1-i\nu_2)
\Big)\delta_{\partial U} \notag\\
&\quad
-\frac1\pi g(w)\Delta \1_U
+\frac{\lambda}{\pi}\Delta\delta_0. \label{eq:boundary-complex}
\end{align}

Equivalently, using
\[
\bar w(\nu_1+i\nu_2)=w\cdot \nu+i\,w\cdot \tau,
\qquad
w(\nu_1-i\nu_2)=w\cdot \nu-i\,w\cdot \tau,
\]
the first-order boundary term may be written as
\begin{equation}\label{eq:boundary-normal-tangent}
-2e^{-\pi|w|^2}
\Big(
(3-\pi|w|^2)(w\cdot \nu)
+i(\pi|w|^2-1)(w\cdot \tau)
\Big)\delta_{\partial U}.
\end{equation}

Finally, since
\[
\nabla g(w)=-\pi(2-\pi|w|^2)e^{-\pi|w|^2}w,
\]
one has
\[
\partial_\nu g
=
-\pi(2-\pi|w|^2)e^{-\pi|w|^2}(w\cdot \nu),
\]
and therefore
\[
-\frac1\pi g\,\Delta \1_U
=
(2-\pi|w|^2)e^{-\pi|w|^2}(w\cdot \nu)\,\delta_{\partial U}
-\frac1\pi g\,\partial_\nu\delta_{\partial U}
\]
in the standard shorthand
\[
\langle \partial_\nu\delta_{\partial U},\varphi\rangle
:=
-\int_{\partial U}\partial_\nu\varphi\,ds.
\]
Hence the full boundary contribution in \(\Delta u\) can be expressed as
\begin{align}
&e^{-\pi|w|^2}
\Big(
-4(w\cdot \nu)-2i(\pi|w|^2-1)(w\cdot \tau)
\Big)\delta_{\partial U} \notag\\
&\qquad
-\frac1\pi(1-\pi|w|^2)e^{-\pi|w|^2}\partial_\nu\delta_{\partial U}. \label{eq:boundary-final}
\end{align}
\end{lemma}

\begin{proof}
The formulas for \(\pw \1_U\) and \(\pwb \1_U\) are the standard distributional derivatives of a
characteristic function with \(C^1\) boundary. The identity for \(\Delta \1_U\) follows from
\[
\langle \Delta \1_U,\varphi\rangle
=
\langle \1_U,\Delta \varphi\rangle
=
\int_U \Delta \varphi\,dA
=
\int_{\partial U}\partial_\nu\varphi\,ds
\]
by the divergence theorem.

Substituting the formulas for \(\pw \1_U\) and \(\pwb \1_U\) into
\eqref{eq:second-order-full} gives \eqref{eq:boundary-complex}. The normal/tangential rewriting
follows from the identities
\[
\bar w(\nu_1+i\nu_2)=w\cdot \nu+i\,w\cdot \tau,
\qquad
w(\nu_1-i\nu_2)=w\cdot \nu-i\,w\cdot \tau.
\]

For the derivative boundary term, one computes
\[
\langle -\tfrac1\pi g\,\Delta \1_U,\varphi\rangle
=
-\frac1\pi\langle \Delta \1_U,g\varphi\rangle
=
-\frac1\pi\int_{\partial U}\partial_\nu(g\varphi)\,ds
\]
and expands
\[
\partial_\nu(g\varphi)=(\partial_\nu g)\varphi+g\,\partial_\nu\varphi.
\]
Using the formula for \(\nabla g\) yields the stated decomposition. Combining that with
\eqref{eq:boundary-normal-tangent} gives \eqref{eq:boundary-final}.
\end{proof}

\begin{remark}\label{rem:real-part}
The boundary expressions above are complex-valued. In the later rigidity discussion, one must
therefore be careful not to use implicitly that the compactly supported solution \(u\) of
\eqref{eq:first-order-u} is real-valued: there is no reason to expect that in general.

If we set
\[
v:=\Re u,
\]
then, since the Laplacian has real coefficients,
\[
\Delta v=\Re(\Delta u).
\]
Taking the real part of \eqref{eq:boundary-final} gives
\begin{align}
\Delta v
&=
4(1-\pi|w|^2)e^{-\pi|w|^2}\1_U
-4e^{-\pi|w|^2}(w\cdot \nu)\,\delta_{\partial U} \notag\\
&\quad
-\frac1\pi(1-\pi|w|^2)e^{-\pi|w|^2}\partial_\nu\delta_{\partial U}
+\frac{\lambda}{\pi}\Delta\delta_0. \label{eq:real-part-laplacian}
\end{align}
Equivalently, using \(g(w)=(1-\pi|w|^2)e^{-\pi|w|^2}\),
\begin{equation}\label{eq:real-part-laplacian-g}
\Delta v
=
4g(w)\1_U
-2(3-\pi|w|^2)e^{-\pi|w|^2}(w\cdot \nu)\,\delta_{\partial U}
-\frac1\pi g(w)\Delta \1_U
+\frac{\lambda}{\pi}\Delta\delta_0.
\end{equation}
Thus passing to the real part removes the tangential boundary term
\[
-2i(\pi|w|^2-1)e^{-\pi|w|^2}(w\cdot \tau)\,\delta_{\partial U},
\]
and retains only the normal boundary contribution.
\end{remark}

\section{Boundary-value interpretation}

\begin{proposition}\label{prop:bvp}
Assume \(\partial U\) is of class \(C^1\). Let \(u\) be a compactly supported distribution
satisfying
\[
\pw u
=
\bar w\,e^{-\pi|w|^2}\1_U
-\frac1\pi(1-\pi|w|^2)e^{-\pi|w|^2}\pw \1_U
+\frac{\lambda}{\pi}\pw\delta_0.
\]
Assume moreover that \(u\) is represented by a \(C^1\)-function on \(\overline U\setminus\{0\}\),
that \(u\equiv 0\) on the exterior side of \(\partial U\), and that \(0\in U\). Then:
\begin{enumerate}[label=\textnormal{(\roman*)}]
\item the boundary trace of \(u\) is
\begin{equation}\label{eq:dirichlet-trace}
u(w)
=
-\frac1\pi(1-\pi|w|^2)e^{-\pi|w|^2}
\qquad (w\in \partial U);
\end{equation}
\item in \(U\setminus\{0\}\),
\begin{equation}\label{eq:interior-first-order}
\pw u=\bar w\,e^{-\pi|w|^2};
\end{equation}
hence
\begin{equation}\label{eq:interior-laplacian}
\Delta u=4(1-\pi|w|^2)e^{-\pi|w|^2}
\qquad \text{in }U\setminus\{0\};
\end{equation}
\item on \(\partial U\), if \(\nu\) denotes the outward unit normal and \(\tau\) the positively
oriented unit tangent, then
\begin{equation}\label{eq:normal-tangent-derivatives}
\partial_\nu u
=
2(w\cdot \nu)e^{-\pi|w|^2},
\qquad
\partial_\tau u
=
2(w\cdot \tau)e^{-\pi|w|^2};
\end{equation}
\item consequently, the boundary must satisfy the compatibility condition
\begin{equation}\label{eq:compatibility}
(1-\pi|w|^2)(w\cdot \tau)=0
\qquad (w\in \partial U).
\end{equation}
\end{enumerate}
\end{proposition}

\begin{proof}
Since \(u\equiv 0\) on the exterior side, the distribution \(\pw u\) may be written near
\(\partial U\) as
\[
\pw u
=
\1_U\,\pw u + u\,\pw \1_U,
\]
where \(u\) in the second term denotes the interior trace on \(\partial U\).
Comparing this with the given equation and using the fact that \(0\notin \partial U\), so the
term \(\pw\delta_0\) is irrelevant near the boundary, we obtain the coefficient identity
\[
u\,\pw \1_U
=
-\frac1\pi(1-\pi|w|^2)e^{-\pi|w|^2}\pw \1_U.
\]
This gives \eqref{eq:dirichlet-trace}.

Away from \(\partial U\cup\{0\}\), the terms \(\pw \1_U\) and \(\pw\delta_0\) vanish, so the
equation reduces to \eqref{eq:interior-first-order}. Applying \(4\pwb\) yields
\eqref{eq:interior-laplacian}.

To compute the boundary derivatives, recall that
\[
2\pw
=
(\nu_1-i\nu_2)(\partial_\nu-i\partial_\tau).
\]
Using \eqref{eq:interior-first-order} on \(\partial U\) from the interior side gives
\[
(\nu_1-i\nu_2)(\partial_\nu-i\partial_\tau)u
=
2\bar w\,e^{-\pi|w|^2}.
\]
Multiplying by \(\nu_1+i\nu_2\), we obtain
\[
\partial_\nu u-i\partial_\tau u
=
2(\nu_1+i\nu_2)\bar w\,e^{-\pi|w|^2}.
\]
Now
\[
(\nu_1+i\nu_2)\bar w=w\cdot \nu+i\,w\cdot \tau,
\]
so taking real and imaginary parts yields \eqref{eq:normal-tangent-derivatives}.

Finally, since \(u\in C^1(\overline U\setminus\{0\})\), the tangential derivative of the
boundary trace must coincide with the tangential derivative computed from the interior. Differentiating
\eqref{eq:dirichlet-trace} along \(\tau\), we get
\[
\partial_\tau u
=
2(2-\pi|w|^2)e^{-\pi|w|^2}(w\cdot \tau).
\]
Comparing with \eqref{eq:normal-tangent-derivatives} gives
\[
2(2-\pi|w|^2)e^{-\pi|w|^2}(w\cdot \tau)
=
2e^{-\pi|w|^2}(w\cdot \tau),
\]
which is exactly \eqref{eq:compatibility}.
\end{proof}

\section{Why the previous Laplacian argument fails}

The earlier version of the note tried to continue from a global Poisson equation of the form
\[
\Delta u
=
\text{bulk term in }U
+\text{a first-order boundary term}
+\text{a point source at }0.
\]
That structure is what made the subsequent steps possible:
\begin{enumerate}[label=\textnormal{(\arabic*)}]
\item identify exterior vanishing of \(u\);
\item compare the coefficient of a single boundary distribution to recover an explicit
formula for \(\pwb u\) on \(\partial U\);
\item subtract the singular part and build an anti-holomorphic quantity \(q\).
\end{enumerate}

However, the corrected differentiated equation \eqref{eq:second-order-full} has a different
structure:
\begin{enumerate}[label=\textnormal{(\arabic*)}]
\item a bulk term in \(U\);
\item a term involving \(\pwb \1_U\);
\item a term involving \(\pw \1_U\);
\item a second-order boundary distribution \(\pwb\pw \1_U\);
\item a point-source term \(\Delta\delta_0\).
\end{enumerate}

In particular, one no longer has a single first-order boundary distribution whose coefficient
can be read off as \(\pwb u\) on \(\partial U\). The appearance of the second-order boundary
term \(\pwb\pw \1_U\) breaks the coefficient-comparison argument used in the previous version.

\begin{remark}
Therefore the present note does \emph{not} establish the disk theorem. The corrected
Paley--Wiener / Ehrenpreis--Malgrange step yields a valid compact-support equation, but it is
a first-order equation. Rebuilding the rigidity argument from \eqref{eq:first-order-u}, or
finding a different reduction that produces a clean Poisson problem, remains open at the
level of this write-up.
\end{remark}

\section{How analyticity would close the argument}

The compatibility condition from Proposition~\ref{prop:bvp},
\begin{equation}\label{eq:compatibility-repeat}
(1-\pi|w|^2)(w\cdot \tau)=0
\qquad (w\in \partial U),
\end{equation}
is already very close to a rigidity statement. The remaining issue is propagation along the
boundary.

\begin{remark}\label{rem:analytic-closure}
Suppose one could justify, from the corrected first-order formulation, the hypotheses needed to
apply the Kinderlehrer--Nirenberg theorem near every boundary point \(p\in \partial U\) with
\(p\neq 0\). Then \(\partial U\) would be real-analytic, and the compatibility condition
\eqref{eq:compatibility-repeat} would force \(\partial U\) to be a circle centered at the
origin.
\end{remark}

\begin{proof}
Assume \(\partial U\) is real-analytic and connected. Let
\[
\gamma:I\to \partial U
\]
be a real-analytic local parametrization by arclength, so \(\tau(s)=\gamma'(s)\). Define
\[
R(s):=|\gamma(s)|^2.
\]
Then \(R\) is real-analytic and
\[
R'(s)=2\,\gamma(s)\cdot \tau(s).
\]
Hence the compatibility condition \eqref{eq:compatibility-repeat} becomes
\[
(1-\pi R(s))R'(s)=0
\qquad (s\in I).
\]

Now consider the real-analytic function
\[
F(s):=(1-\pi R(s))R'(s).
\]
Since \(F\equiv 0\), analytic continuation shows that on each connected chart one of the following
must happen:
\begin{enumerate}[label=\textnormal{(\roman*)}]
\item \(R(s)\equiv \pi^{-1}\), in which case the boundary arc lies on the circle
\(|w|=\pi^{-1/2}\);
\item \(R'(s)\equiv 0\), in which case \(R\) is constant and the boundary arc lies on a circle
centered at the origin.
\end{enumerate}
Thus every analytic boundary arc is contained in a circle centered at the origin.

Because \(\partial U\) is connected and real-analytic, two neighboring analytic arcs must agree
on their common endpoint with matching tangent. Two centered circles meeting with the same tangent
must have the same radius. Hence all local arcs have the same radius, and therefore
\[
\partial U=\{\,|w|=R\,\}
\]
for some \(R>0\).

If \(U\) is simply connected, this implies that \(U\) is the disk of radius \(R\) centered at the
origin.
\end{proof}

The point of the previous remark is that the boundary compatibility condition is not the real
obstruction anymore. Once analyticity of the free boundary is available, the geometry propagates
automatically. The remaining missing ingredient is therefore the analytic free-boundary input
itself: one must show that the corrected first-order equation furnishes the hypotheses needed to
invoke Kinderlehrer--Nirenberg near \(\partial U\).

\section{Summary}

For the \((h_1,h_1,U)\) problem, the justified conclusions are:
\begin{enumerate}[label=\textnormal{(\arabic*)}]
\item the Bargmann eigenvalue identity \eqref{eq:bargmann-h1};
\item the equivalent conjugated identity \eqref{eq:bargmann-h1-conjugated};
\item the compactly supported first-order distribution \(\omega\) in \eqref{eq:def-omega};
\item the vanishing relation \(\langle \omega,e^{\pi wz}\rangle=0\);
\item the corresponding compactly supported first-order equation \(\pw u=\omega\);
\item the differentiated second-order equation \eqref{eq:second-order-full}.
\end{enumerate}

What is \emph{not} currently justified is the passage from these facts to the Laplacian-based
free-boundary argument that would force \(U\) to be a disk.

\end{document}
